% =====================
\section{作業計画}

\subsection{作業分解構造WBSと管理番号}

作業分解構造（Work Breakdown Structure：WBS）を図\ref{fig:wbs}に示す．
WBSは設計フェーズ（100番台），実装フェーズ（200番台），テストフェーズ（300番台）の
3フェーズで構成され，計60の活動タスクからなる（\texttt{base\_plan/wbs\_v2.xlsx}参照）．
各活動タスクは8/80ルール（8時間以上80時間未満）に準拠して設定した．

\begin{figure}[H]
  \centering
  \fbox{\begin{minipage}{0.92\linewidth}
    \centering\vspace{4cm}
    \textit{（WBS図を後から挿入：\texttt{base\_plan/wbs\_v2.drawio}）}
    \vspace{4cm}
  \end{minipage}}
  \caption{作業分解構造WBS}
  \label{fig:wbs}
\end{figure}

\subsection{アローダイアグラムとクリティカルパス}

（後から記載・図挿入予定）

\begin{figure}[H]
  \centering
  \fbox{\begin{minipage}{0.92\linewidth}
    \centering\vspace{4cm}
    \textit{（アローダイアグラムを後から挿入）}
    \vspace{4cm}
  \end{minipage}}
  \caption{アローダイアグラムとクリティカルパス}
  \label{fig:arrow}
\end{figure}

\subsection{役割分担}

各担当の役割と担当WBSタスクを表\ref{tab:roles}に示す．

\begin{table}[H]
  \centering
  \begin{tabularx}{\linewidth}{llXl}
    \toprule
    担当 & メンバー & 主な担当タスク & WBS番号 \\
    \midrule
    A & 長見東磨・飛田 & 指揮者人形機構設計・サーボ/LED制御・回路配線 & 120系・210〜216系 \\
    B & 長野志哉・慶田 & センサー配線・観測SW・状態機械実装 & 131〜132系・221〜225系・310〜320系 \\
    C & 長山魁良・慶田 & 楽譜配列・演奏制御SW・通信プロトコル実装 & 141〜142系・231〜234系・330系 \\
    D & 中村海惺・尾添 & Processing音響合成・エフェクト実装 & 151〜152系・241〜242系 \\
    全員 & — & FBS/PBS/WBS作成・MOE確定・システムテスト・発表準備 & 111〜114系・340〜343系 \\
    \bottomrule
  \end{tabularx}
  \caption{役割分担}
  \label{tab:roles}
\end{table}

\subsection{ガントチャート}

（後から記載・図挿入予定）

\begin{figure}[H]
  \centering
  \fbox{\begin{minipage}{0.92\linewidth}
    \centering\vspace{5cm}
    \textit{（ガントチャートを後から挿入）}
    \vspace{5cm}
  \end{minipage}}
  \caption{ガントチャート}
  \label{fig:gantt}
\end{figure}
